% generated by scripts/paper_conclusion.py, rule conclusion-v1, from reports/final/final_primary.json; do not edit
We built a procedural benchmark of drawn directed graphs in which every label and every optimal repair is checked by an independent verifier, answer-flipping twins bound any image-blind predictor at $0.50$ per task, and splits are grouped by root graph with held-out styles, families and sizes. On it we ran a pre-registered, exposure-matched comparison of repair training against question-answering and structural fine-tuning, with oracle-graph and parse-and-solve conditions that separate perception errors from reasoning errors. We release the generator, verifier, protocol, frozen evaluation code and per-sample results so the comparison can be rerun or extended.

Repair training (arm~D) did worse than the reference fine-tuning arm on the pre-registered six-cell statistic; the hypothesis is not supported at this scale. On the locked test the difference D$-$B is $-0.105$ (95\% interval [$-0.117$, $-0.093$], mean of 3 seeds). Within this setting, learning to make a minimal verified edit did not carry over to non-edit questions as well as training on those questions directly; Section~\ref{sec:negative} records the development observations that pointed the same way. Per-split, per-task and secondary comparisons are in Section~\ref{sec:test-results}.

This is a pilot-scale study: all adapters were trained on $512$ synthetic root graphs with one 2B base model (Qwen3.5-2B, LoRA), and nothing here measures real diagrams. A negative result at pilot scale does not rule out transfer with more data, a larger model or a different edit budget; the released protocol can be rerun unchanged to test each of these. Section~\ref{sec:limitations} lists the controls we did not run and the other limits of this study.
